\documentclass[10pt]{extarticle}
\usepackage[a4paper,margin=1.5cm]{geometry}
\usepackage[T1]{fontenc}
\usepackage{amsmath,amssymb}
\usepackage{graphicx}
\usepackage[export]{adjustbox}
\usepackage{multicol}
\usepackage{enumitem}
\usepackage{booktabs,array,tabularx}
\setlength{\columnsep}{0.7cm}
\setlength{\parindent}{0pt}
\setlength{\parskip}{3pt}
\pagestyle{empty}
\makeatletter
% figure registry (figdefs.tex is generated by the extract step)
\newcommand{\qfigdef}[4]{\@namedef{qf@file@#1}{#2}\@namedef{qf@w@#1}{#3}\@namedef{qf@wide@#1}{#4}}
\newcommand{\qf@check}[1]{\@ifundefined{qf@file@#1}{\PackageError{score2}{Unknown figure id #1}{}}{}}
\newcommand{\qf@img}[4][]{\includegraphics[width=#2\dimexpr\@nameuse{qf@w@#4}pt\relax,max width=#3,max height=0.45\textheight,#1]{\@nameuse{qf@file@#4}}}
% Full-width questions: finalize wraps every question that contains a wide figure in \qfullbegin...\qfullend,
% so the WHOLE question (stem, figures, options) is set across the page in reading order, then the two columns
% resume. The question's head - its number, stem and first figure (\qfullsplit marks the end of the head) -
% is kept on one page, so a number is never left at a page foot with its figure overleaf; the rest may break
% across pages normally, so no big gaps are left. Without a \qfullsplit, the whole question is kept together
% if it is under 45% of a page. Never write these yourself.
\newif\ifqf@full
\newif\ifqf@boxed
\newsavebox{\qf@box}
\newcommand{\qf@open}{\global\qf@fulltrue\par\global\qf@boxedtrue
  \setbox\qf@box\vbox\bgroup\hsize\textwidth\linewidth\textwidth}
\newcommand{\qf@emit}[1]{\par\egroup\global\qf@boxedfalse
  \ifdim\dimexpr\ht\qf@box+\dp\qf@box\relax<#1\textheight\noindent\usebox{\qf@box}\par\else\unvbox\qf@box\fi}
\newcommand{\qfullsplit}{\ifqf@boxed\qf@emit{0.85}\fi}
\newcommand{\qf@close}{\ifqf@boxed\qf@emit{0.45}\fi\global\qf@fullfalse}
\newcommand{\qfullbegin}{\end{multicols}\qf@open}
\newcommand{\qfullend}{\qf@close\begin{multicols}{2}}
\newcommand{\qfullstart}{\qf@open}
\newcommand{\qfulldone}{\qf@close}
% \qfig[scale]{id}: block figure, centred, directly where it is written (size set by the tool).
% Fallback if a wide figure is not inside a full-width question: break out of the columns in place - no floats.
\newcommand{\qfig}[2][1]{\qf@check{#2}%
  \ifqf@full\qf@blk{#1}{#2}\else
  \ifnum\@nameuse{qf@wide@#2}=1 \expandafter\@firstoftwo\else\expandafter\@secondoftwo\fi
  {\end{multicols}\begin{center}\qf@img{#1}{\textwidth}{#2}\end{center}\begin{multicols}{2}}%
  {\qf@blk{#1}{#2}}\fi}
\newcommand{\qf@blk}[2]{\par\smallskip{\centering\qf@img{#1}{\linewidth}{#2}\par}\smallskip}
% \qfigi[scale]{id}: inline figure (inside option lists / tables / minipages); never breaks columns.
\newcommand{\qfigi}[2][1]{\qf@check{#2}\qf@img[valign=c]{#1}{\linewidth}{#2}}
\makeatother
% structure macros - lint reads these to check completeness and figure placement
\newcommand{\papersec}[1]{\par\bigskip{\centering\Large\bfseries PAPER-#1\par}\medskip}
\newcommand{\partsec}[2]{\par\medskip\noindent{\large\bfseries PART-#1: #2}\par\smallskip}
\newcommand{\sect}[1]{\par\smallskip\noindent\textbf{#1}\par}
\newcommand{\qn}[1]{\par\smallskip\noindent\textbf{#1.}~}
\newcommand{\paraq}[2]{\par\smallskip\noindent\textbf{Paragraph for Questions #1 and #2}\par}
\newenvironment{qoptions}{\begin{itemize}[leftmargin=1.8em,itemsep=1pt,topsep=2pt,label={}]}{\end{itemize}}
\newcommand{\opt}[2]{\item[(#1)] #2}
% figure registry (generated by score2_tools.py finalize)
\qfigdef{p1_phy_s0_q01_1}{SRGBT - Test 02 (01 Aug 2021)-img/p1_phy_s0_q01_1.png}{93.9}{0}
\qfigdef{p1_phy_s0_q02_1}{SRGBT - Test 02 (01 Aug 2021)-img/p1_phy_s0_q02_1.png}{118.3}{0}
\qfigdef{p1_phy_s0_q03_1}{SRGBT - Test 02 (01 Aug 2021)-img/p1_phy_s0_q03_1.png}{174.1}{0}
\qfigdef{p1_phy_s0_q04_1}{SRGBT - Test 02 (01 Aug 2021)-img/p1_phy_s0_q04_1.png}{86.1}{0}
\qfigdef{p1_phy_s0_q05_1}{SRGBT - Test 02 (01 Aug 2021)-img/p1_phy_s0_q05_1.png}{76.1}{0}
\qfigdef{p1_phy_s0_q06_1}{SRGBT - Test 02 (01 Aug 2021)-img/p1_phy_s0_q06_1.png}{54.2}{0}
\qfigdef{p1_phy_s0_q07_1}{SRGBT - Test 02 (01 Aug 2021)-img/p1_phy_s0_q07_1.png}{67.5}{0}
\qfigdef{p1_phy_s0_q08_1}{SRGBT - Test 02 (01 Aug 2021)-img/p1_phy_s0_q08_1.png}{106.3}{0}
\qfigdef{p1_phy_s0_q08_2}{SRGBT - Test 02 (01 Aug 2021)-img/p1_phy_s0_q08_2.png}{163.7}{0}
\qfigdef{p1_phy_s0_q10_1}{SRGBT - Test 02 (01 Aug 2021)-img/p1_phy_s0_q10_1.png}{69.1}{0}
\qfigdef{p1_phy_s0_q12_1}{SRGBT - Test 02 (01 Aug 2021)-img/p1_phy_s0_q12_1.png}{131.0}{0}
\qfigdef{p1_phy_s0_q15_1}{SRGBT - Test 02 (01 Aug 2021)-img/p1_phy_s0_q15_1.png}{85.2}{0}
\qfigdef{p1_phy_s0_q16_1}{SRGBT - Test 02 (01 Aug 2021)-img/p1_phy_s0_q16_1.png}{113.6}{0}
\qfigdef{p1_phy_s0_q17_1}{SRGBT - Test 02 (01 Aug 2021)-img/p1_phy_s0_q17_1.png}{109.2}{0}
\qfigdef{p1_phy_s0_q18_1}{SRGBT - Test 02 (01 Aug 2021)-img/p1_phy_s0_q18_1.png}{95.2}{0}
\qfigdef{p1_phy_s0_q20_1}{SRGBT - Test 02 (01 Aug 2021)-img/p1_phy_s0_q20_1.png}{81.4}{0}
\qfigdef{p1_chem_s0_q21_1}{SRGBT - Test 02 (01 Aug 2021)-img/p1_chem_s0_q21_1.png}{245.2}{0}
\qfigdef{p1_chem_s0_q23_1}{SRGBT - Test 02 (01 Aug 2021)-img/p1_chem_s0_q23_1.png}{320.7}{1}
\qfigdef{p1_chem_s0_q28_1}{SRGBT - Test 02 (01 Aug 2021)-img/p1_chem_s0_q28_1.png}{182.2}{0}
\qfigdef{p1_chem_s0_q32_4}{SRGBT - Test 02 (01 Aug 2021)-img/p1_chem_s0_q32_4.png}{273.3}{1}
\qfigdef{p1_chem_s0_q35_1}{SRGBT - Test 02 (01 Aug 2021)-img/p1_chem_s0_q35_1.png}{177.0}{0}
\qfigdef{p1_chem_s0_q36_1}{SRGBT - Test 02 (01 Aug 2021)-img/p1_chem_s0_q36_1.png}{260.4}{1}
\raggedcolumns
\makeatletter\@beginparpenalty=10000\makeatother
\newcolumntype{L}[1]{>{\raggedright\arraybackslash}p{#1}}
\begin{document}
\begin{center}{\Large\bfseries SRGBT --- Test 02 (01 Aug 2021)}\end{center}
\begin{multicols}{2}
\raggedcolumns
\papersec{1}
\partsec{1}{PHYSICS}
\qn{1} The system in figure is released from rest from the position shown. After blocks have moved distance H/3, collar B is removed and block A \& C continue to move. The speed of C just before it strikes the ground is\par
\qfig{p1_phy_s0_q01_1}
\begin{qoptions}
\opt{A}{$\dfrac43\sqrt{\mathrm{gH}}$}
\opt{B}{$2\sqrt{\dfrac{\mathrm{gH}}3}$}
\opt{C}{$\dfrac{\sqrt{13\mathrm{gH}}}3$}
\opt{D}{$2\sqrt{2\mathrm{gH}}$}
\end{qoptions}
\qn{2} Three blocks A, B, C of mass M, M $-$ m$_0$, m$_0$ respectively consists a system as shown in figure. The string and pulley are massless and all surfaces are frictionless. When system is released the block C moves downwards against the vertical surface of block B. At any instant, the velocity of block A is 2 m/s. Then the magnitude of velocity of block C w.r.t ground at considered instant is\par
\qfig{p1_phy_s0_q02_1}
\begin{qoptions}
\opt{A}{3 m/s}
\opt{B}{$\sqrt{20}$ m/s}
\opt{C}{2 m/s}
\opt{D}{$\sqrt5$ m/s}
\end{qoptions}
\qn{3} In the setup shown, blocks of masses m, 3m and 4m are placed on a frictionless horizontal surface and the free end `P' of the thread is being pulled by a constant force 2F. Find the acceleration of the free end P.\par
\qfig{p1_phy_s0_q03_1}
\begin{qoptions}
\opt{A}{8 F/m}
\opt{B}{4 F/m}
\opt{C}{16 F/m}
\opt{D}{None of these}
\end{qoptions}
\qn{4} A uniform chain is at rest partially on the inclined and partially hanging vertically. Coefficient of friction between chain and inclined is $\mu=\dfrac1{2\sqrt3}$. The ratio of $\dfrac{\mathrm L_{\max}}{\mathrm L_{\min}}$ is (L$_{\max}$ = maximum length of chain which kept on inclined so that chain remains at rest, L$_{\min}$ = minimum length of chain which kept on inclined so that chain remains at rest)\par
\qfig{p1_phy_s0_q04_1}
\begin{qoptions}
\opt{A}{$\dfrac74$}
\opt{B}{$\dfrac65$}
\opt{C}{$\dfrac75$}
\opt{D}{None of these}
\end{qoptions}
\qn{5} The following figure shows an accelerating conveyor belt inclined at an angle 37$^\circ$ above horizontal. The coefficient of friction between the belt and block is `1'. The least time in which block can reach the top, starting from rest at the bottom is :\par
\qfig{p1_phy_s0_q05_1}
\begin{qoptions}
\opt{A}{2s}
\opt{B}{4s}
\opt{C}{1s}
\opt{D}{0.5s}
\end{qoptions}
\qn{6} In the system shown in the diagram all surfaces are smooth, pulley and strings are ideal. If $\vec a_\mathrm A$ and $\vec a_\mathrm B$ are the accelerations of the two blocks, then just after the system is released from rest, then choose correct statement(s) :\par
\qfig{p1_phy_s0_q06_1}
\begin{qoptions}
\opt{A}{$|\vec a_\mathrm A|=|\vec a_\mathrm B|$}
\opt{B}{$\vec a_\mathrm A\perp\vec a_\mathrm B$}
\opt{C}{Acceleration of A relative to B is vertically downwards.}
\opt{D}{Normal force exerted by B on A is zero.}
\end{qoptions}
\qn{7} A wedge of mass m$_1$ and a block of mass m$_2$ are in equilibrium as shown. Inclined surface of the wedge has an inclination $\alpha$ with the horizontal. Each surface is frictionless. The normal reaction on the wedge may be\par
\qfig{p1_phy_s0_q07_1}
\begin{qoptions}
\opt{A}{m$_2$g cos $\alpha$}
\opt{B}{m$_2$ g sin $\alpha$ cos $\alpha$}
\opt{C}{m$_1$g + m$_2$g cos$^2\alpha$}
\opt{D}{m$_1$g + m$_2$g sin $\alpha$ cos $\alpha$}
\end{qoptions}
\qn{8} A small metal ball of mass m is being pulled gradually on a fixed frictionless hemisphere as shown in figure. Radii of ball and that of the pulley are much smaller than that of the hemisphere. Then the values of F (magnitude of pulling) force and R (contact force between the ball and the hemisphere change) ?\par
\qfig{p1_phy_s0_q08_1}
\begin{qoptions}
\opt{A}{F = mg$\dfrac{\ell}{\mathrm H}$}
\opt{B}{F = mg$\dfrac{\mathrm r}{\mathrm H}$}
\opt{C}{R = mg$\dfrac{\ell}{\mathrm H}$}
\opt{D}{R = mg$\dfrac{\mathrm r}{\mathrm H}$}
\end{qoptions}
\paraq{9}{10}Three blocks A, B and C are connected with the help of pulleys, string and springs (all are ideal) as shown in the figure. All the surfaces are smooth. Take g = 10 m/s$^2$ (Masses of blocks A, B \& C are 10 kg, 5 kg \& 20$\sqrt3$ kg respectively)\par
\qfig{p1_phy_s0_q08_2}
\qn{9} The tension in the string joining B and C is
\begin{qoptions}
\opt{A}{10 $\sqrt3$ N}
\opt{B}{100 $\sqrt3$ N}
\opt{C}{20 N}
\opt{D}{50 N}
\end{qoptions}
\qn{10} If the string between B and C is cut, then the magnitude of acceleration of the block B becomes
\begin{qoptions}
\opt{A}{10 $\sqrt3$ m/s$^2$}
\opt{B}{$\dfrac{20}{3\sqrt3}$ m/s$^2$}
\opt{C}{20 $\sqrt3$ m/s$^2$}
\opt{D}{None of these}
\end{qoptions}
\paraq{11}{12}4 identical particles are connected with four massless springs and placed at corners of square as shown in figure. Strings (i) and (ii) are massless and vertical and whole system is in equilibrium. Then choose correct statement.\par
\qfig{p1_phy_s0_q10_1}
\qn{11} Just after string 1 is cut, then acceleration of A will be
\begin{qoptions}
\opt{A}{2g downward}
\opt{B}{zero}
\opt{C}{g upward}
\opt{D}{g$\sqrt2$ at 45$^\circ$ with AD}
\end{qoptions}
\qn{12} Just after spring between C \& D is cut only, acceleration of C will be-
\begin{qoptions}
\opt{A}{g down wards}
\opt{B}{zero}
\opt{C}{2g downward}
\opt{D}{g$\sqrt2$ at 45$^\circ$ with BC}
\end{qoptions}
\paraq{13}{14}A horizontal conveyor belt is running at a constant speed v$_\mathrm b$ = 3.0 m/s. A small disc enters the belt moving horizontally with a velocity v$_0$ = 4.0 m/s that is perpendicular to velocity of the belt. Coefficient of friction between the disc and the belt is 0.50. [Take g = 10 m/s$^2$]\par
\qfig{p1_phy_s0_q12_1}
\qn{13} What should be the minimum width of the belt be so that the disc always remain on the belt ?
\begin{qoptions}
\opt{A}{0.9 m}
\opt{B}{1.6 m}
\opt{C}{2.0}
\opt{D}{2.5 m}
\end{qoptions}
\qn{14} What will be the time at which there will be minimum speed of the disc relative to ground ?
\begin{qoptions}
\opt{A}{0.5 sec}
\opt{B}{1 sec}
\opt{C}{16/25 sec}
\opt{D}{4/3 sec}
\end{qoptions}
\qn{15} A smooth tube in the shape of a rectangle with rounded corners is placed in a vertical plane, as shown in figure. Two ball bearings are introduced at the upper right-hand corner. One travels by path AB and the other by path CD to arrive at the lower left-hand corner. If the time taken by them is t$_\mathrm{AB}$ and t$_\mathrm{CD}$ respectively, find the value of t$_\mathrm{AB}$/t$_\mathrm{CD}$.\par
\qfig{p1_phy_s0_q15_1}
\qn{16} Two hemispheres of negligible of radius R combined together by two minimum forces `F' to form a sphere (completely filled with water) as shown in figure. A small pin hole of radius a at top of the sphere as shown in figure. If the system is in equilibrium at $\tan\theta=\dfrac x2$. Find the value of x. (gravity is present and no other support is given to hemispheres)\par
\qfig{p1_phy_s0_q16_1}
\qn{17} A unifrom chain of mass M = 10 kg hangs in vertical plane as shown is figure. Find mass of segment AP of the chain. [Take g = 10 m/s$^2$]\par
\qfig{p1_phy_s0_q17_1}
\qn{18} A side view of a simplified form of vertical latch B is as shown. The lower member A can be pushed forward in its horizontal channel. The sides of the channels are smooth, but at the interfaces of A and B, which are at 45$^\circ$ with the horizontal, there exists a static coefficient of friction $\mu$ = 0.4. If the minimum force F = 2P(in N) that must be applied horizontally to A to start motion of the latch B if it has a mass m = 0.6 kg. Find the value of P\par
\qfig{p1_phy_s0_q18_1}
\qn{19} Consider a book lying on the floor and a string with length $l$ = 2.0m whose one end is attached to the centre of the upper side of book's cover. The other end of the string is held in height h = 1.2m above the floor. Find the coefficient of friction between the book and the floor knowing that the book which is of mass m = 1 kg is dragged at constant speed with force of magnitude F = 10N. Assume that the book does not open and neglect height of book. If your answer is x fill value of 3x. (g = 10 m/s$^2$)
\qn{20} In the figure is a string attached at one end to wall and other end to block B of mass 1kg. The total length of string is 35 cm but the vertical part is of length 10 cm. B is resting against vertical wall of A of mass 3 kg free to move on the floor. The minimum value of coefficient of static friction between A and floor to avoid slipping anywhere is $\dfrac4{\mathrm k}$. Find k. [Neglect size of pulley]\par
\qfig{p1_phy_s0_q20_1}
\raggedcolumns
\papersec{1}
\partsec{2}{CHEMISTRY}
\qn{21} Two flasks A and B of equal volume containing NH$_3$ and HCl gases, are connected by a narrow tube of negligible volume. The two gases were prevented from mixing by stopper fitted in connecting tube. For further detail of experiment refer to the given figure. What will be final pressure in each flask when passage connecting two tubes are opened. Assume ideal gas behaviour of NH$_3$ and HCl gas and the reaction.\par
NH$_3$(g) + HCl(g) $\longrightarrow$ NH$_4$Cl(s) is quantitative.\par
\qfig{p1_chem_s0_q21_1}
\begin{qoptions}
\opt{A}{40 mm Hg}
\opt{B}{60 mm Hg}
\opt{C}{20 mm Hg}
\opt{D}{10 mm Hg}
\end{qoptions}
\qn{22} Consider the following orbitals in an atom.\par
S$_1$ represents set of orbital(s) which is/are spherically symmetrical with two radial nodes.\par
S$_2$ represents set of orbital(s) which is/are double dumb--bell or dumble with ring \& have no radial nodes.\par
S$_3$ represents set of orbital(s) with zero angular momentum \& three radial nodes.\par
S$_4$ represents set of orbital(s) having one planar \& one radial node.\par
Using the above information, identify which of the following option is incorrect.
\begin{qoptions}
\opt{A}{Maximum degeneracy from the above set if the orbitals belong to He is 5.}
\opt{B}{Number of orbitals from the above set which have $|$m$|\leq$ l = 8.}
\opt{C}{Maximum number of orbital whose ``n + $\ell$'' is same but angular momentum differs can be 5.}
\opt{D}{Orbitals S$_1$ and S$_3$ will have same radial probability distribution graph.}
\end{qoptions}
\qfullbegin
\qn{23} A closed vessel of known volume containing known amount of ideal gaseous substance `A' was observed for variation of pressure with temperature. The expected graph was to be like as in (i). However, the actual observations revealed the graph to be like (ii). The deviation was attributed to polymerization of gas molecules as nA(g) $\rightleftharpoons$ A$_\mathrm n$(g) . If it is known that the above reaction gives only 50\% yield.\par
Find the value of n to which the gas A is being polymerized.\par
\qfig{p1_chem_s0_q23_1}
\qfullsplit
\begin{qoptions}
\opt{A}{2}
\opt{B}{3}
\opt{C}{4}
\opt{D}{5}
\end{qoptions}
\qfullend
\qn{24} Which of the following option is true w.r.t wave mechanical model.
\begin{qoptions}
\opt{A}{For a 1s orbital of He$^+$, radial probability density attains a maxima at $\dfrac{\mathrm a_0}2$ where a$_0$ is Bohr radius of hydrogen.}
\opt{B}{For the same value of ``n'', as ``$\ell$'' increases the value of ``r$_\mathrm{maximum}$'' (radial distance where total radical probability is maximised) increases.}
\opt{C}{In case of H-atom degeneracy of 3$^\mathrm{rd}$ excited state is 16.}
\opt{D}{For 3d orbital there exists no space where probability of finding e$^-$ is 0.}
\end{qoptions}
\qn{25} Two flasks A and B have equal volume. `A' is maintained at 300 K and `B' at 600 K. `A' contains H$_2$ gas while `B' has on equal mass of CH$_4$ gas. Assuming ideal behaviour for both the gases identify the correct statements.\par
(I) Greater number of molecules are present in flask `B'.\par
(II) High pressure is attained in flask `A'.\par
(III) Compressibility factor of both flasks is same.\par
(IV) Total kinetic energy is same in both the flasks.\par
(V) Molar kinetic energy is greater for flaks `B'.\par
(VI) Molecules are travelling with the same speed in both flasks.
\begin{qoptions}
\opt{A}{I, II, III}
\opt{B}{I, III, V}
\opt{C}{II, III, V}
\opt{D}{II, III, IV}
\end{qoptions}
\qn{26} In a hydrogen like sample electron is in 2$^\mathrm{nd}$ excited state, the binding energy of 4$^\mathrm{th}$ state of this sample is 7.65 eV then:
\begin{qoptions}
\opt{A}{A 15 eV photon can set free the electron from the second excited state of this sample.}
\opt{B}{Three different types of photons will be observed if electrons make transition upto ground state from second excited state.}
\opt{C}{If 13 eV photon is used then K.E. of the ejected electron is 1 eV.}
\opt{D}{2$^\mathrm{nd}$ line of Balmer series of this sample has same energy value as 1$^\mathrm{st}$ excitation energy of H--atoms.}
\end{qoptions}
\qn{27} For a real gas obeying vander Waal's equation which of the following are true ?
\begin{qoptions}
\opt{A}{Internal energy of gas is dependent on temperature only.}
\opt{B}{If Z $>$ 1 then forces of repulsion are dominant.}
\opt{C}{At high pressure Z = $\left(1+\dfrac{\mathrm{Pb}}{\mathrm{RT}}\right)$}
\opt{D}{Work done in expansion of a real gas is more than for an ideal gas under isothermal conditions.}
\end{qoptions}
\qn{28} Which of the following statements is/are correct ?
\begin{qoptions}
\opt{A}{The energy of electron in an orbital in the hydrogen atoms depends upon the value of principal quantum number only.}
\opt{B}{The energy of an electron is largely determined by its principal quantum number}
\opt{C}{The principal quantum number is a measure of the most probable distance of finding the electron around the nucleus.}
\opt{D}{The value of principal quantum number of 24$^\mathrm{th}$ electron in an atom is 3.}
\end{qoptions}
\paraq{29}{30}The figure shows \textbf{initial conditions} of a uniform cylinder with frictionless pistons A and B held in shown position by mechanical stoppers.\par
\qfig{p1_chem_s0_q28_1}
If the mechanical stoppers holding piston A and B as shown in figure are removed and in the mean time N$_2$O$_4$ and O$_3$ gases undergoes following reaction completely.\par
[Assume that temperature remains constant.]\par
N$_2$O$_4$ $\longrightarrow$ 2NO$_2$\par
2O$_3$ $\longrightarrow$ 3O$_2$
\qn{29} Which statement is incorrect after attaining final equilibrium state
\begin{qoptions}
\opt{A}{Position of piston A is 45 cm from left end}
\opt{B}{Position of piston B is 60 cm from left end}
\opt{C}{Total moles of gases in all compartments is 18}
\opt{D}{Pistion `B' final \& initial position are different}
\end{qoptions}
\qn{30} Which statement is correct after all the piston have attained their final positions\par
(assume $\sigma$ to be same for all gases)
\begin{qoptions}
\opt{A}{Z$_{11}$ will be highest for molecules of gas in compartment 1}
\opt{B}{Mean free path will be longest for molecules of gas in compartment 2}
\opt{C}{Z$_{11}$ will be same for all compartments}
\opt{D}{Mean free path will be same for molecules of gas in all compartments}
\end{qoptions}
\paraq{31}{32}Both de--Broglie hypothesis and Heisenberg uncertainty principle establish that microscope particles have more prominent wave nature.\par
Heisenberg realised that, in the world of very small particles one cannot measure any property of a particle without interacting with it in some way. This introduces an unavoidable uncertainty into the result. One can never measure all the properties exactly.\par
Minimum uncertainty in position of any body is given by wavelength of light.\par
\textbf{Heisenberg's uncertainty principle :}\par
$\Delta\mathrm x\;.\;\Delta\mathrm P\;\geq\;\dfrac{\mathrm h}{4\pi}=\dfrac{\hbar}2$\par
(uncertainty in position, uncertainty in momentum)\par
\textbf{De--Broglie hypothesis :}\par
$\lambda=\dfrac{\mathrm h}{\mathrm P}$\par
where $\lambda$ is de--Broglie wavelength of any body having non relativistic momentum P.
\qn{31} To image a virus of linear size d, through what voltage electrons (mass `m' and charge `e') should be accelerated in electron microscope.
\begin{qoptions}
\opt{A}{V $\geq\dfrac{2\pi^2\hbar^2}{\mathrm{med}^2}$}
\opt{B}{V $\leq\dfrac{2\pi^2\hbar^2}{\mathrm{med}^2}$}
\opt{C}{V $\geq\dfrac{\mathrm h}{2\mathrm{med}^2}$}
\opt{D}{V $\leq\dfrac{\hbar}{2\mathrm{med}^2}$}
\end{qoptions}
\qn{32} Starting from uncertainty in position and momentum in Heisenbergs uncertainty principle, what will be the correct relation of uncertainty in de--Broglie wavelength and position of a particle in one dimension ?
\begin{qoptions}
\opt{A}{$\Delta\lambda.\Delta\mathrm x\geq\dfrac{\hbar}2$}
\opt{B}{$\Delta\lambda.\Delta\mathrm x\geq\dfrac{\mathrm h}{2\pi}$}
\opt{C}{$\Delta\lambda.\Delta\mathrm x\geq-\dfrac{\lambda^2}{4\pi}$}
\opt{D}{$\Delta\lambda.\Delta\mathrm x\geq\dfrac{\lambda^2}{2\pi}$}
\end{qoptions}
\qfullbegin
\paraq{33}{34}Two containers A and B of capacity 1 litre and 2 litre respectively are connected by tube of negligible volume. The tube is initially closed by stopcock and `A' contains small amount of H$_2$O($l$) and `B' contain initially pure COCl$_2$(g) at pressure of 100 torr. The gas COCl$_2$ partially dissociates at experimental condition into CO(g) and Cl$_2$(g). Now the stopcock connecting two container is opened and sufficient time is given to attain final equilibrium. The final pressure at equilibrium is found to be 100 torr in both containers.\par
Using given data and taking following assumption calculate the mole fraction of CO(g) in final state.\par
\textbf{Given : Vapour pressure of H$_2$O at experimental temperature `T' is equal to 20 torr.}\par
\textbf{Assumption-1 :} Volume of H$_2$O($l$) is small and can be neglected but sufficient to create vapour pressure in both container.\par
\textbf{Assumption-2 :} COCl$_2$, CO and Cl$_2$ gas do not dissolve in H$_2$O($l$).\par
\textbf{Assumption-3 :} Whole experiment is carried out under isothermal condition.\par
\qfig{p1_chem_s0_q32_4}
\qfullend
\qn{33} The ratio of number of moles of CO(g) to total number of moles of gases in the final state is:
\begin{qoptions}
\opt{A}{$\dfrac4{15}$}
\opt{B}{$\dfrac8{15}$}
\opt{C}{$\dfrac6{15}$}
\opt{D}{$\dfrac2{15}$}
\end{qoptions}
\qn{34} Ratio of the partial pressure of CO(g) and COCl$_2$(g) in final state is:
\begin{qoptions}
\opt{A}{5 : 1}
\opt{B}{4 : 1}
\opt{C}{1 : 4}
\opt{D}{1 : 5}
\end{qoptions}
\qn{35} A container is divided into two parts in the ratio 1 : 3 by a fixed semi permeable membrane (SPM) (which permits movement of NO$_2$ only) as shown. At temperature `T' kelvin chamber I consists of undissociated N$_2$O$_4$ at a pressure of 20 mm and chamber II contains small amount of liquid H$_2$O. The chambers are now heated to the temperature `1.2 T' kelvin at which N$_2$O$_4$ dissociates by 50\%.\par
Given : The vapour pressure of water at `T' kelvin is 20 mm of Hg and the vapour pressure at `1.2T' kelvin is 30 mm of Hg.\par
\qfig{p1_chem_s0_q35_1}
Assume H$_2$O does not completely evaporate at `T' or `1.2 T' Kelvin.\par
Calculate abcd where\par
ab = two digit number representing pressure (in mm) in chamber I at `1.2 T' Kelvin.\par
cd = two digit number representing pressure (in mm) in chamber II at `1.2 T' Kelvin.
\qn{36} Using the following information, calculate the value of a four digit number ``abcd''.\par
``a'' = ratio of shortest wavelength of Balmer series to that of longest wavelength in U--V region of atomic spectrum of hydrogen.\par
``b'' = Stopping potential in volts when EMR of 124 nm strikes a metal surface having critical wavlength equal to 310 nm.\par
``c'' = Number of possible exchanges of electron in d--orbital which can occur in a d$^4$ configuration.\par
``d'' = Number of de--Broglie waves in the ``N'' orbit of hydrogen atom.
\qfullbegin
\qn{37} \qfig{p1_chem_s0_q36_1}
\qfullsplit
The above apparatus consists of three temperature jacketed 12.315 litre bulbs connected by stop cocks. Bulbs `A' consist of NH$_{3(g)}$, CO$_{2(g)}$ \& N$_{2(g)}$ at 27$^\circ$C and a total pressure of 60 atm. Bulb `B' and `C' are empty and maintained at temperature of $-$73$^\circ$C and $-$173$^\circ$C respectively. Initially stop-cocks are closed and volume of the connecting lines is zero.\par
\textbf{Given :}\par
(i) NH$_{3(g)}$ liquifies at $-$20$^\circ$C, CO$_{2(g)}$ liquifies at $-$80$^\circ$C and N$_{2(g)}$ liquifies at $-$180$^\circ$C.\par
(ii) Vapour pressure of NH$_{3(l)}$ at $-$73$^\circ$C = 1 atm\par
(iii) Vapour pressure of NH$_{3(l)}$ at $-$173$^\circ$C is negligible\par
(iv) Vapour pressure of CO$_{2(l)}$ is negligible.\par
(v) On opening stop-cock (i) only, the pressure in containers becomes 17 atm and, on opening both stop-cocks pressure decreases to 4 atm.\par
Calculate \textbf{`abcd'}\par
where ab = moles of NH$_3$ taken\par
cd = moles of nitrogen taken\par
\textbf{[Give your answer as ab + cd and fill OMR sheet]}
\qfullend
\qn{38} If velocity of electron moving in n$^\mathrm{th}$ Bohr's orbit of H--atom is `v$_\mathrm n$', what is the velocity of a hypothetical particle of charge same as electron and mass two times to that of electron revolving in the same orbit.\par
Assume the nucleus to be non stationary. [Given : 1840 m$_\mathrm e$ = m$_\mathrm p$]\par
\textbf{[Express your answer in terms of v}$'_\mathrm n$ = v$_\mathrm n\times\dfrac{(1840+\mathrm x)}{1842}$\textbf{]. Then calculate x.}
\qn{39} How many correct statements matchings (where P$_\mathrm C$ = Critical pressure \& T$_\mathrm C$ = Critical temeprature).\par
{\footnotesize\setlength{\tabcolsep}{1.5pt}\noindent\begin{tabular}{|c|L{2.6cm}|c|L{3.2cm}|}
\hline
(i) & Real Gas below critical temperature but Pressure = 0.2 times P$_\mathrm C$ & $\Rightarrow$ & Gas can be liquefied by applying higher pressure without changing its T \\
\hline
(ii) & Super critical CO$_2$ at Pressure = 20 times P$_\mathrm C$ & $\Rightarrow$ & Gas always obey Boyle's and Charle's laws at these conditions. \\
\hline
(iii) & Real gas above critical temperature but pressure = 50 times P$_\mathrm C$ & $\Rightarrow$ & Attractive forces are definitely dominant over repulsive forces at these conditions. \\
\hline
(iv) & Real gas at T = 0.5 T$_\mathrm C$ but pressure = 0.1 times P$_\mathrm C$ & $\Rightarrow$ & Gas can be liquefied by changing temperature keeping pressure constant \\
\hline
\end{tabular}}\par
\qn{40} An unknown orbital is having the following equation for its wavefunction.\par
$\psi(\mathrm r,\theta,\phi)=\dfrac23\left(\dfrac{\mathrm Z}{3\mathrm a_0}\right)^{3/2}(\sigma-1)(\sigma^2-8\sigma+12)\sigma\mathrm e^{-\sigma/2}\cos\theta$\par
where $\sigma=\dfrac{2\mathrm{Zr}}{\mathrm a_0}$, a$_0$ is constant = 0.529 \AA\ and all other parameters have usual meaning. Based on this data answer the question that follow.\par
What is the number of total nodes of the orbital ? (Total nodes = raidal nodes + angular nodes)
\raggedcolumns
\papersec{1}
\partsec{3}{MATHEMATICS}
\qn{41} If a, b, c, d $\in$ R then the equation (x$^2$ + ax $-$ 3b) (x$^2$ $-$ cx + b) (x$^2$ $-$ dx + 2b) = 0 has
\begin{qoptions}
\opt{A}{6 real roots}
\opt{B}{3 real roots}
\opt{C}{4 real roots}
\opt{D}{at least 2 real roots}
\end{qoptions}
\qn{42} The number of values of $\lambda$ for which [x$^2$ $-$ ($\lambda$ $-$ 2)x + $\lambda^2$] [x$^2$ + $\lambda$x + (2$\lambda$ $-$ 1)] is a perfect square is
\begin{qoptions}
\opt{A}{1}
\opt{B}{2}
\opt{C}{0}
\opt{D}{None of these}
\end{qoptions}
\qn{43} If 11 A.M.s are inserted between 28 and 10, then number of integral A.M's is
\begin{qoptions}
\opt{A}{5}
\opt{B}{6}
\opt{C}{7}
\opt{D}{8}
\end{qoptions}
\qn{44} Let x, y, z be positive real numbers such that xyz $\geq\left(\dfrac{\mathrm x}3+\dfrac{\mathrm y}3+\dfrac{\mathrm z}3\right)^3$ and 3x + 4y + 5z = 12, then the value of x$^3$ + y$^4$ + z$^5$ is equal to \underline{\hspace{1.5cm}}
\begin{qoptions}
\opt{A}{3}
\opt{B}{4}
\opt{C}{5}
\opt{D}{6}
\end{qoptions}
\qn{45} For real a, b, c ; a + b + c = 2, a$^2$ + b$^2$ + c$^2$ = 6 and a$^3$ + b$^3$ + c$^3$ = 8, then $\dfrac{(1-\mathrm a)(1-\mathrm b)(1-\mathrm c)}5$ is equal to \underline{\hspace{1.5cm}}.
\begin{qoptions}
\opt{A}{$-$1}
\opt{B}{1}
\opt{C}{2}
\opt{D}{0}
\end{qoptions}
\qn{46} If a, b, c be three unequal positive quantities in H.P. then
\begin{qoptions}
\opt{A}{a$^{100}$ + c$^{100}$ $>$ 2b$^{100}$}
\opt{B}{a$^3$ + c$^3$ $>$ 2b$^3$}
\opt{C}{a$^5$ + c$^5$ $>$ 2b$^5$}
\opt{D}{a$^2$ + c$^2$ $>$ 2b$^2$}
\end{qoptions}
\qn{47} Let a$_1$, a$_2$, ......, a$_n$ represent real numbers such that $\displaystyle\sum_{i=1}^{n}\mathrm a_i=1$ and a$_i\geq0\;\forall$ i $\in$ \{1, 2, ......, n\}. If S = $\displaystyle\sum_{i=1}^{n-1}\mathrm a_i\mathrm a_{i+1}$ , then S can take which of the following values
\begin{qoptions}
\opt{A}{0}
\opt{B}{$\dfrac14$}
\opt{C}{$\dfrac23$}
\opt{D}{1}
\end{qoptions}
\qn{48} If the roots of the equation x$^2$ $-$ px $-$ 1 = 0 and x$^2$ $-$ qx $-$ 1 = 0 form (in a suitable order) an arithmetic progression with four members then
\begin{qoptions}
\opt{A}{p = $\dfrac2{\sqrt3}$, q = $\dfrac{-2}{\sqrt3}$}
\opt{B}{p = $\dfrac4{\sqrt3}$, q = $\dfrac{-4}{\sqrt3}$}
\opt{C}{p = $\dfrac{-4}{\sqrt3}$, q = $\dfrac4{\sqrt3}$}
\opt{D}{p = $\dfrac{-2}{\sqrt3}$, q = $\dfrac2{\sqrt3}$}
\end{qoptions}
\paraq{49}{50}a, b, c, d is an A.P. with distinct integer terms and d = a$^2$ + b$^2$ + c$^2$.
\qn{49} The common difference is
\begin{qoptions}
\opt{A}{1}
\opt{B}{2}
\opt{C}{3}
\opt{D}{None of these}
\end{qoptions}
\qn{50} The value of a + 2b + 3c + 4d is
\begin{qoptions}
\opt{A}{0}
\opt{B}{10}
\opt{C}{100}
\opt{D}{None of these}
\end{qoptions}
\paraq{51}{52}Consider the quadratic equation ax$^2$ $-$ (2a $-$ 1)x + a $-$ 3 = 0 and answer the following questions :
\qn{51} If both the roots are real and are of opposite sign, then complete set of values of a is
\begin{qoptions}
\opt{A}{(0, 3)}
\opt{B}{($-\infty$, 0) $\cup$ (3, $\infty$)}
\opt{C}{$\left\{\dfrac12\right\}$}
\opt{D}{None of these}
\end{qoptions}
\qn{52} If a is given as $\dfrac{\mathrm n^2+\mathrm n}2$, where `n' is rational number, then both the roots are necessarily
\begin{qoptions}
\opt{A}{Integers}
\opt{B}{rational numbers}
\opt{C}{even integers}
\opt{D}{none of these}
\end{qoptions}
\paraq{53}{54}If roots of the equation x$^4$ $-$ 12x$^3$ + bx$^2$ + cx + 81 = 0 are positive, then
\qn{53} Value of b is
\begin{qoptions}
\opt{A}{$-$54}
\opt{B}{54}
\opt{C}{27}
\opt{D}{$-$27}
\end{qoptions}
\qn{54} Value of c is
\begin{qoptions}
\opt{A}{108}
\opt{B}{$-$108}
\opt{C}{54}
\opt{D}{$-$54}
\end{qoptions}
\qn{55} f(x) = x$^6$ $-$ 2x$^3$ $-$ 8 and g(x) = x$^2$ + 2x + 4. Let $\alpha_1$, $\alpha_2$, ..............., $\alpha_6$ are the roots of equation f(x) = 0, find g($\alpha_1$)$\cdot$ g($\alpha_2$) ...... g($\alpha_6$).
\qn{56} The number of real roots of (6 $-$ x)$^4$ + (8 $-$ x)$^4$ = 16 is
\qn{57} If $\alpha+\dfrac1\alpha$ \& $2-\beta-\dfrac1\beta$ ($\alpha$, $\beta>0$) are the roots of the quadratic equation x$^2$ $-$ 2(a + 1)x + a $-$ 3 = 0, then the sum of all integral value of `a'.
\qn{58} If the sum of roots of the quadratic equation ($-$a + 1)x$^2$ + (a$^2$ $-$ a + 4)x $-$ 2a + 3 = 0 is minimum, then the value of a(a $>$ 1) is
\qn{59} If the polynomial f(x) = 4x$^4$ $-$ ax$^3$ + bx$^2$ $-$ cx + 5 where a, b, c $\in$ R has four positive real roots say r$_1$, r$_2$ r$_3$ and r$_4$, such that $\dfrac{\mathrm r_1}2+\dfrac{\mathrm r_2}4+\dfrac{\mathrm r_3}5+\dfrac{\mathrm r_4}8=1$ . Find the value of (a $-$ 10).
\qn{60} Given a$^2$ + b$^2$ + c$^2$ $\leq$ 16 and S = 3a + 4b + 10c where a, b, c $\in$ R. If the maximum value of S is $\sqrt{\mathrm N}$ , then the value of N is equal to
\end{multicols}
\end{document}
